\documentclass[11pt]{article}
\usepackage[margin=1in]{geometry}
\usepackage{amsmath, amssymb, amsthm}
\usepackage{mathtools}
\usepackage{xcolor}
\usepackage{hyperref}
\hypersetup{colorlinks=true, linkcolor=blue, urlcolor=blue}

\title{From the harmonic toy to the sech$^2$ disk:\\
A complete derivation of the $\sigma_w$/$\sigma_z$ mismatch asymmetry}
\author{}
\date{}

\begin{document}
\maketitle

\section{The question}

We have a 1D vertical phase-space simulation. Stars are sampled at $t=0$ from a Gaussian
\begin{equation}
f_0(z_0,w_0)=\frac{1}{2\pi\sigma_z\sigma_w}\exp\!\left[-\frac{z_0^2}{2\sigma_z^2}-\frac{w_0^2}{2\sigma_w^2}\right]
\label{eq:f0}
\end{equation}
and integrated forward to a time $t$ in a vertical potential $\Phi(z;h)$, where $h$ is the
dark-matter scale-height parameter we want to fit. We observe
$\{(z_{\rm obs}^{(i)},w_{\rm obs}^{(i)})\}_{i=1}^{N}$ at $t$ and recover $h$ by maximum
likelihood: for a trial $h$, back-integrate each star to $t=0$, then evaluate the
log-prior of the back-integrated cloud under \emph{assumed} sigmas
$(\sigma_{z,\rm fit},\sigma_{w,\rm fit})$.

The empirical observation we want to explain:
\begin{itemize}
\item Shrink $\sigma_{w,\rm fit}$ by a factor $r$ ($r=10$ or $100$). The fitted $h$ runs
away from the truth.
\item Shrink $\sigma_{z,\rm fit}$ by the same factor $r$. The fit barely moves.
\end{itemize}
We will derive the reason in two steps. First, the harmonic potential, where every
quantity is in closed form. Then the real (sech$^2$ + DM) potential, where the same
structure holds but the algebra has to be done as orbit averages.

% =========================================================
\section{Harmonic potential: full closed-form derivation}
% =========================================================

\subsection{The potential, the Hamiltonian, and the equations of motion}

Consider a single particle of unit mass in
\begin{equation}
\Phi(z) = \tfrac12\nu^2 z^2,
\qquad
H(z,w) = \tfrac12 w^2 + \tfrac12\nu^2 z^2.
\end{equation}
Hamilton's equations are
\begin{equation}
\dot z = \frac{\partial H}{\partial w} = w,
\qquad
\dot w = -\frac{\partial H}{\partial z} = -\nu^2 z.
\label{eq:eom-harm}
\end{equation}
The frequency $\nu$ is the analog of the dark-matter parameter $h$ in this toy: it is
the parameter the fitter tries to recover.

\subsection{The flow}

The general solution of \eqref{eq:eom-harm} starting from $(z_0,w_0)$ at $t=0$ is
\begin{equation}
\boxed{\;
z(t) = z_0\cos(\nu t) + \frac{w_0}{\nu}\sin(\nu t),\qquad
w(t) = -\nu z_0\sin(\nu t) + w_0\cos(\nu t).
\;}
\label{eq:flow-harm}
\end{equation}
This is a rotation by angle $\theta=\nu t$ in the rescaled coordinates
$(z, w/\nu)$ — symplectic and area-preserving.

\subsection{The back-integration map}

Inverting \eqref{eq:flow-harm} at trial frequency $\nu$, given observations
$(z_{\rm obs},w_{\rm obs})$ at time $t$:
\begin{equation}
\boxed{\;
z_0(\nu) = z_{\rm obs}\cos(\nu t) - \frac{w_{\rm obs}}{\nu}\sin(\nu t),\qquad
w_0(\nu) = \nu z_{\rm obs}\sin(\nu t) + w_{\rm obs}\cos(\nu t).
\;}
\label{eq:back-harm}
\end{equation}
Throughout the rest of this section, $(z_{\rm obs},w_{\rm obs})$ are
\emph{fixed by the data}, generated by \eqref{eq:flow-harm} at the true frequency
$\nu_t$, while $\nu$ is the trial frequency the fitter is varying.

\subsection{The truth flow inserted into the back-integration}

\paragraph{What we want and why.}
The back-map \eqref{eq:back-harm} gives $z_0(\nu)$ and $w_0(\nu)$ as linear functions of
the \emph{observations} $z_{\rm obs}, w_{\rm obs}$. That's not yet useful for computing
variances, because the observations are themselves random (they depend on the unknown
truth $z_0, w_0$). We want $z_0(\nu)$ and $w_0(\nu)$ expressed in terms of the
\emph{truth} $z_0, w_0$ — the variables whose statistics we know (Gaussian with
$\sigma_z, \sigma_w$).

\paragraph{The idea: chain two linear maps.}
\begin{itemize}
\item The back-map (eq.~\eqref{eq:back-harm}) is linear in $(z_{\rm obs},w_{\rm obs})$.
\item The forward truth-flow (eq.~\eqref{eq:flow-harm} at $\nu_t$) makes
$(z_{\rm obs},w_{\rm obs})$ linear in $(z_0,w_0)$.
\item Substituting the second into the first, $(z_0(\nu),w_0(\nu))$ becomes linear in
$(z_0,w_0)$.
\end{itemize}
The coefficients of this composite linear map are exactly $A,B,C,D$. Concretely:
$A(\nu)$ and $C(\nu)$ are the coefficients of $z_0$ and $w_0$ in $z_0(\nu)$;
$B(\nu)$ and $D(\nu)$ are the same for $w_0(\nu)$. No tricks — substitute and collect.

\paragraph{Explicit substitution.}
Substituting the truth flow
\begin{equation}
z_{\rm obs} = z_0\cos\theta_t + \frac{w_0}{\nu_t}\sin\theta_t,
\qquad
w_{\rm obs} = -\nu_t z_0\sin\theta_t + w_0\cos\theta_t,
\qquad \theta_t\equiv\nu_t t,
\end{equation}
into \eqref{eq:back-harm} (and writing $\theta\equiv\nu t$) gives
\begin{equation}
z_0(\nu) = A(\nu)\,z_0 + C(\nu)\,w_0,
\qquad
w_0(\nu) = B(\nu)\,z_0 + D(\nu)\,w_0,
\label{eq:zw-back-decomp}
\end{equation}
with
\begin{align}
A(\nu) &= \cos\theta_t\cos\theta + \tfrac{\nu_t}{\nu}\sin\theta_t\sin\theta,\\[2pt]
C(\nu) &= \tfrac{1}{\nu_t}\sin\theta_t\cos\theta - \tfrac{1}{\nu}\cos\theta_t\sin\theta,\\[2pt]
B(\nu) &= \nu\cos\theta_t\sin\theta - \nu_t\sin\theta_t\cos\theta,\\[2pt]
D(\nu) &= \tfrac{\nu}{\nu_t}\sin\theta_t\sin\theta + \cos\theta_t\cos\theta.
\end{align}
Sanity check at $\nu=\nu_t$: $A=D=1$ and $B=C=0$, so $z_0(\nu_t)=z_0$ and $w_0(\nu_t)=w_0$.
The back-integration recovers the truth.

\subsection{The variances of the back-integrated cloud}

Because $\langle z_0 w_0\rangle=0$ under the prior \eqref{eq:f0}:
\begin{equation}
\boxed{\;
V_z(\nu) \equiv \langle z_0(\nu)^2\rangle = \sigma_z^2 A(\nu)^2 + \sigma_w^2 C(\nu)^2,
\;}
\end{equation}
\begin{equation}
\boxed{\;
V_w(\nu) \equiv \langle w_0(\nu)^2\rangle = \sigma_z^2 B(\nu)^2 + \sigma_w^2 D(\nu)^2.
\;}
\end{equation}
At truth: $V_z(\nu_t)=\sigma_z^2$, $V_w(\nu_t)=\sigma_w^2$. The fit at $\nu_t$ recovers the
correct sigmas.

\subsection{The likelihood}

Up to $\nu$-independent constants, the per-star negative log-likelihood the fitter
optimises is
\begin{equation}
\mathrm{NLL}(\nu) = \frac{V_z(\nu)}{2\sigma_{z,\rm fit}^2}+\frac{V_w(\nu)}{2\sigma_{w,\rm fit}^2}.
\label{eq:NLL-harm}
\end{equation}
Differentiating in $\nu$:
\begin{equation}
\boxed{\;
\mathrm{NLL}'(\nu) = 0
\quad\Longleftrightarrow\quad
\frac{V_z'(\nu)}{\sigma_{z,\rm fit}^2} + \frac{V_w'(\nu)}{\sigma_{w,\rm fit}^2} = 0.
\;}
\label{eq:gradbal}
\end{equation}
This is the equation the fitter solves for $\nu_\star$ (the MLE).

\subsection{First derivatives at $\nu=\nu_t$}

Differentiate $A,B,C,D$ once with respect to $\nu$ and evaluate at $\nu=\nu_t$. Keeping
all terms (not yet asymptotic in $t$):
\begin{align}
A'(\nu_t) &= -\frac{\sin^2\theta_t}{\nu_t},\\[2pt]
C'(\nu_t) &= -\frac{t}{\nu_t} + \frac{\sin 2\theta_t}{2\nu_t^2},\\[2pt]
B'(\nu_t) &= \nu_t t + \frac{\sin 2\theta_t}{2},\\[2pt]
D'(\nu_t) &= +\frac{\sin^2\theta_t}{\nu_t}.
\end{align}
Then, using $A(\nu_t)=D(\nu_t)=1$ and $B(\nu_t)=C(\nu_t)=0$:
\begin{equation}
V_z'(\nu_t) = 2\sigma_z^2 A(\nu_t)A'(\nu_t) + 2\sigma_w^2 C(\nu_t)C'(\nu_t)
= -\frac{2\sigma_z^2\sin^2\theta_t}{\nu_t},
\end{equation}
\begin{equation}
V_w'(\nu_t) = 2\sigma_z^2 B(\nu_t)B'(\nu_t) + 2\sigma_w^2 D(\nu_t)D'(\nu_t)
= +\frac{2\sigma_w^2\sin^2\theta_t}{\nu_t}.
\end{equation}
\begin{equation}
\boxed{\;
V_z'(\nu_t) = -\frac{2\sigma_z^2\sin^2\theta_t}{\nu_t},
\qquad
V_w'(\nu_t) = +\frac{2\sigma_w^2\sin^2\theta_t}{\nu_t}.
\;}
\label{eq:firstderiv}
\end{equation}

\paragraph{Truth-balance check.}
With $\sigma_{z,\rm fit}=\sigma_z$ and $\sigma_{w,\rm fit}=\sigma_w$,
\eqref{eq:gradbal} at $\nu_t$ reads
\begin{equation}
-\frac{2\sin^2\theta_t}{\nu_t} + \frac{2\sin^2\theta_t}{\nu_t} = 0,
\end{equation}
so the MLE sits exactly at $\nu_t$. The two variance gradients pull in opposite
directions, and at truth sigmas they cancel exactly.

\subsection{Second derivatives at $\nu=\nu_t$ — the key asymmetry}

We need
\begin{equation}
V_z''(\nu_t) = 2\sigma_z^2[(A')^2 + AA'']\big|_{\nu_t}
            + 2\sigma_w^2[(C')^2 + CC'']\big|_{\nu_t},
\end{equation}
\begin{equation}
V_w''(\nu_t) = 2\sigma_z^2[(B')^2 + BB'']\big|_{\nu_t}
            + 2\sigma_w^2[(D')^2 + DD'']\big|_{\nu_t}.
\end{equation}
We need $A'',\,B'',\,C'',\,D''$ at $\nu_t$. Direct differentiation, keeping only the
leading $t^2$ terms (large-$t$, many-orbits regime), gives
\begin{equation}
A''(\nu_t)\simeq -t^2,
\quad
D''(\nu_t)\simeq -t^2,
\end{equation}
and the leading large-$t$ behaviours
\begin{equation}
C'(\nu_t)\simeq -\frac{t}{\nu_t}\Rightarrow [C'(\nu_t)]^2 \simeq \frac{t^2}{\nu_t^2},
\qquad
B'(\nu_t)\simeq \nu_t t \Rightarrow [B'(\nu_t)]^2\simeq \nu_t^2 t^2.
\end{equation}
Plugging in, with $A=D=1$ and $B=C=0$ at $\nu_t$, only the $A\,A''$, $D\,D''$,
$(B')^2$, $(C')^2$ terms survive at order $t^2$:
\begin{align}
V_z''(\nu_t) &\simeq 2\sigma_z^2(-t^2) + 2\sigma_w^2\,\frac{t^2}{\nu_t^2}
= 2t^2\!\left(\frac{\sigma_w^2}{\nu_t^2} - \sigma_z^2\right),\\[4pt]
V_w''(\nu_t) &\simeq 2\sigma_z^2\,\nu_t^2 t^2 + 2\sigma_w^2(-t^2)
= 2t^2(\sigma_z^2\nu_t^2 - \sigma_w^2).
\end{align}

\subsection{The dimensionless ratio $R$}

Define
\begin{equation}
\boxed{\;
R \equiv \frac{\sigma_w}{\nu_t\sigma_z}
= \frac{\text{velocity-dispersion-induced amplitude}}{\text{position-dispersion amplitude}}.
\;}
\end{equation}
In the simulation $\sigma_w\sim 20\;\mathrm{km/s}$, $\sigma_z\sim 10\;\mathrm{pc}$,
$\nu_t\sim 0.04\;\mathrm{Myr^{-1}}$ (after unit-conversion), so $R\sim 50$.
We will assume $R\gg 1$.

The second derivatives become
\begin{equation}
V_z''(\nu_t) = 2\sigma_z^2 t^2(R^2-1) \;\xrightarrow{R\gg 1}\; +\,2\sigma_z^2 t^2 R^2,
\end{equation}
\begin{equation}
V_w''(\nu_t) = 2\sigma_w^2 t^2(1/R^2 - 1) \;\xrightarrow{R\gg 1}\; -\,2\sigma_w^2 t^2.
\end{equation}
Note $V_w''<0$: $V_w(\nu)$ has a local \emph{maximum} at $\nu_t$ (the variance of the
back-integrated $w$-coordinate \emph{decreases} when you bump $\nu$). $V_z(\nu)$ has a
local \emph{minimum}. The two contributions to $\mathrm{NLL}''$ have opposite signs;
their truth-sigma sum
\begin{equation}
\mathrm{NLL}''(\nu_t)\big|_{\rm truth}
= \frac{V_z''(\nu_t)}{\sigma_z^2} + \frac{V_w''(\nu_t)}{\sigma_w^2}
\simeq 2t^2(R^2-1) + 2t^2(1/R^2-1) > 0
\end{equation}
is still positive (a true minimum), but dominated by the $V_z$-term:
\begin{equation}
\boxed{\;
\frac{V_z''(\nu_t)/\sigma_z^2}{|V_w''(\nu_t)|/\sigma_w^2}\;\simeq\; R^2.
\;}
\label{eq:Rsquared}
\end{equation}
\textbf{The $z$-direction is $R^2$ times stiffer per unit prior weight than the
$w$-direction.}

\paragraph{Physical picture.}
A trial-frequency error $\delta\nu$ rotates the back-integrated cloud by
$\delta\theta = t\,\delta\nu$ in canonical coordinates. The rotation pumps
\emph{velocity into position} with weight $\sigma_w^2/\nu_t^2$ (large, because the cloud
is hot and the orbits are slow) but \emph{position into velocity} with weight
$\sigma_z^2\nu_t^2$ (small). Hence $V_z$ is highly sensitive to $\nu$, $V_w$ much less.

\subsection{The MLE bias under sigma mismatch}

The Newton step around $\nu_t$:
\begin{equation}
\delta\nu_\star \;\simeq\; -\,\frac{\mathrm{NLL}'(\nu_t)}{\mathrm{NLL}''(\nu_t)}.
\label{eq:newton}
\end{equation}

\paragraph{Case A: only $\sigma_w$ mismatched ($\sigma_{w,\rm fit}=\sigma_w/r$, $\sigma_{z,\rm fit}=\sigma_z$).}
\begin{align}
\mathrm{NLL}'(\nu_t) &= \frac{V_z'(\nu_t)}{\sigma_z^2} + r^2\,\frac{V_w'(\nu_t)}{\sigma_w^2}
= \frac{2\sin^2\theta_t}{\nu_t}\,(r^2 - 1),\\[4pt]
\mathrm{NLL}''(\nu_t) &= \frac{V_z''(\nu_t)}{\sigma_z^2} + r^2\,\frac{V_w''(\nu_t)}{\sigma_w^2}
\simeq 2t^2 R^2 - 2t^2 r^2 \;\xrightarrow{R\gg r}\; 2t^2 R^2.
\end{align}
\begin{equation}
\boxed{\;
\delta\nu_w \;\simeq\; -\,\frac{(r^2-1)\sin^2\theta_t}{\nu_t\,t^2\,R^2}.
\;}
\label{eq:dnuw}
\end{equation}

\paragraph{Case B: only $\sigma_z$ mismatched ($\sigma_{z,\rm fit}=\sigma_z/r$, $\sigma_{w,\rm fit}=\sigma_w$).}
\begin{align}
\mathrm{NLL}'(\nu_t) &= r^2\,\frac{V_z'(\nu_t)}{\sigma_z^2} + \frac{V_w'(\nu_t)}{\sigma_w^2}
= -\,\frac{2\sin^2\theta_t}{\nu_t}\,(r^2 - 1),\\[4pt]
\mathrm{NLL}''(\nu_t) &= r^2\,\frac{V_z''(\nu_t)}{\sigma_z^2} + \frac{V_w''(\nu_t)}{\sigma_w^2}
\simeq 2t^2\,R^2 r^2 - 2t^2 \;\xrightarrow{R^2 r^2 \gg 1}\; 2t^2 R^2 r^2.
\end{align}
\textbf{The numerator is the same magnitude as Case A — the asymmetry is in the
denominator.}
\begin{equation}
\boxed{\;
\delta\nu_z \;\simeq\; +\,\frac{(r^2-1)\sin^2\theta_t}{\nu_t\,t^2\,R^2\,r^2}.
\;}
\label{eq:dnuz}
\end{equation}

\paragraph{The $r^2$ asymmetry.}
\begin{equation}
\boxed{\;
\frac{|\delta\nu_w|}{|\delta\nu_z|} \;=\; r^2.
\;}
\label{eq:rsq-asymmetry}
\end{equation}
At $r=10$: 100$\times$ asymmetry. At $r=100$: $10^4\times$. The $\sigma_w$ mismatch is
overwhelmingly the dominant source of bias.

\subsection{One-sentence intuition}

Mismatching $\sigma_z$ amplifies the curvature in the \emph{already-stiff} $V_z$
direction (which dominates $\mathrm{NLL}''$) — the optimum barely shifts. Mismatching
$\sigma_w$ amplifies the \emph{already-soft} $V_w$ direction (which is sub-dominant) —
the optimum is pushed far before the small $V_w$-curvature can stop it.

% =========================================================
\section{General potential: sech$^2$ disk + DM}
% =========================================================

\subsection{The actual potential}

The simulation uses
\begin{equation}
\Phi(z;h) \;=\; \underbrace{\sum_{i\in\{\text{thin, thick, gas}\}} 4\pi G\rho_i h_i^2 \log\cosh(z/h_i)}_{\Phi_{\rm bary}(z)}
\;+\; \underbrace{4\pi G\rho_{\rm DM}\,h^2\,\log\cosh(z/h)}_{\Phi_{\rm DM}(z;h)}.
\label{eq:realphi}
\end{equation}
The baryon component is fixed; the DM component depends on the fitting parameter $h$.

\subsection{No closed-form flow}

The equations of motion
\begin{equation}
\dot z = w,\qquad \dot w = -\partial_z\Phi(z;h)
\label{eq:eom-real}
\end{equation}
are nonlinear and have \emph{no closed-form solution} for any non-trivial sum of
$\log\cosh$ terms. The forward map
\begin{equation}
(z_0,w_0)\;\longmapsto\;\big(F(z_0,w_0;h,t),\;G(z_0,w_0;h,t)\big),
\end{equation}
where $F=z(t)$ and $G=w(t)$ along the orbit, has to be evaluated numerically. So we
cannot write $V_z(h)$ and $V_w(h)$ as elementary functions. We \emph{can}, however,
compute their derivatives at $h=h_t$ analytically by linearising — the variational
equation.

\subsection{The variational equation}

Define the orbital perturbation
\begin{equation}
\eta(\tau) \;\equiv\; \frac{\partial z(\tau;z_0,w_0,h)}{\partial h}\bigg|_{h=h_t},
\qquad
\dot\eta(\tau) \;=\; \frac{\partial w(\tau;z_0,w_0,h)}{\partial h}\bigg|_{h=h_t},
\end{equation}
i.e.\ the shift in the orbit at time $\tau$ when we bump $h$ slightly, holding
$(z_0,w_0)$ fixed. Differentiating \eqref{eq:eom-real} with respect to $h$ and
evaluating along the true orbit:
\begin{equation}
\boxed{\;
\ddot\eta(\tau) + \Phi_{zz}\!\big(z(\tau);h_t\big)\,\eta(\tau)
\;=\; -\,\Phi_{zh}\!\big(z(\tau);h_t\big),
\qquad \eta(0)=\dot\eta(0)=0,
\;}
\label{eq:vareq}
\end{equation}
where $\Phi_{zz}\equiv\partial^2\Phi/\partial z^2$ and
$\Phi_{zh}\equiv\partial^2\Phi/(\partial z\,\partial h)$.

This is a \emph{linear} ODE along each orbit — a forced harmonic oscillator with a
time-varying ``stiffness'' $\Phi_{zz}(z(\tau);h_t)$ and a known drive
$-\Phi_{zh}(z(\tau);h_t)$. Solvable by the orbit's Green's function (Wronskian of the
homogeneous tangent flow), which itself has to be evaluated numerically.

\subsection{From $\eta(t)$ to $\partial z_0/\partial h$}

Given $(\eta(t),\dot\eta(t))$ along each orbit, the implicit function theorem applied
to $F(z_0(h),w_0(h);h,t)=z_{\rm obs}$ and analogous for $G$ — together with the
symplectic identity $\det J=1$ for $J = \partial(F,G)/\partial(z_0,w_0)$ — gives
\begin{equation}
\frac{\partial z_0}{\partial h}\bigg|_{h_t}
= -\,\eta(t)\,J_{ww} + \dot\eta(t)\,J_{zw},
\end{equation}
\begin{equation}
\frac{\partial w_0}{\partial h}\bigg|_{h_t}
= +\,\eta(t)\,J_{wz} - \dot\eta(t)\,J_{zz},
\end{equation}
where the tangent-flow Jacobian entries
$J_{zz}=\partial F/\partial z_0$, $J_{zw}=\partial F/\partial w_0$,
$J_{wz}=\partial G/\partial z_0$, $J_{ww}=\partial G/\partial w_0$
are themselves solutions of the homogeneous variational equation (same ODE as
\eqref{eq:vareq} but with zero forcing and unit-vector initial conditions).

\subsection{$V_z'$ and $V_w'$ as orbit averages}

Then
\begin{equation}
V_z'(h_t) = 2\Big\langle z_0\,\frac{\partial z_0}{\partial h}\Big\rangle,\qquad
V_w'(h_t) = 2\Big\langle w_0\,\frac{\partial w_0}{\partial h}\Big\rangle,
\end{equation}
and
\begin{equation}
V_z''(h_t)\simeq 2\Big\langle\!\Big(\frac{\partial z_0}{\partial h}\Big)^{\!2}\Big\rangle,
\qquad
V_w''(h_t)\simeq 2\Big\langle\!\Big(\frac{\partial w_0}{\partial h}\Big)^{\!2}\Big\rangle.
\label{eq:Vsecondreal}
\end{equation}
These are exact identities. They cannot be reduced to elementary functions of
$\sigma_z, \sigma_w, h_t$ — they require numerical orbit integration to evaluate. But
the \emph{scaling} is determined by how $\eta(t)$ grows on a typical orbit.

\subsection{Why the harmonic scaling persists}

Action-angle coordinates: in any 1D Hamiltonian, define
\begin{equation}
J(E;h) = \frac{1}{2\pi}\oint w\,dz,
\qquad
\Omega(J;h) = \frac{dE}{dJ},
\end{equation}
so the unperturbed flow is the rigid rotation
$\theta(\tau) = \theta_0 + \Omega(J;h_t)\tau$, with $J$ conserved. On a single orbit
of frequency $\Omega(J)$, the variational equation \eqref{eq:vareq} has the same
linear-secular structure as in the harmonic case: $\eta(\tau)$ grows asymptotically
as
\begin{equation}
\eta(t) \;\sim\; -\frac{t}{\Omega(J)}\,w_0\;+\;\text{oscillatory},
\end{equation}
\begin{equation}
\dot\eta(t) \;\sim\; +\,t\,\Omega(J)\,z_0\;+\;\text{oscillatory},
\end{equation}
where the prefactors come from the orbit average of $\Phi_{zh}$ around the cycle (the
``secular'' part of the perturbation). This is the \emph{exact same parametric
scaling} as the harmonic result, with $\nu\to\Omega(J)$.

Substituting into \eqref{eq:Vsecondreal} and averaging over the action distribution
$\mathcal P(J)$ of the cloud:
\begin{equation}
V_z''(h_t)\;\sim\;2t^2\,\sigma_w^2\,\Big\langle\frac{1}{\Omega(J)^2}\Big\rangle_J,
\qquad
V_w''(h_t)\;\sim\;2t^2\,\sigma_z^2\,\big\langle\Omega(J)^2\big\rangle_J.
\end{equation}
Define the orbit-averaged version of $R$:
\begin{equation}
\boxed{\;
R_{\rm eff}^2 \;\equiv\; \frac{\sigma_w^2\,\langle 1/\Omega^2\rangle}{\sigma_z^2\,\langle\Omega^2\rangle^{-1}}
\;\approx\;\frac{\sigma_w^2}{\sigma_z^2\,\Omega_{\rm eff}^2},
\quad\text{with}\quad \Omega_{\rm eff}\equiv\langle\Omega^{2}\rangle^{1/2}\!\!\!\text{ or similar.}
\;}
\end{equation}
Then
\begin{equation}
\frac{V_z''(h_t)/\sigma_z^2}{|V_w''(h_t)|/\sigma_w^2}\;\simeq\;R_{\rm eff}^2.
\end{equation}
\textbf{All the harmonic results carry through.} Specifically:
\begin{itemize}
\item The first-derivative truth-balance (cancellation at $h=h_t$ with truth sigmas)
holds exactly — it follows from energy conservation and symplecticity, not from
harmonicity.
\item The numerator symmetry of the bias ($\propto r^2-1$ in either case) is exact —
it follows from the form of the Gaussian likelihood alone.
\item The stiffness asymmetry $V_z''/\sigma_z^2 \sim R_{\rm eff}^2\,V_w''/\sigma_w^2$
holds whenever the action-distribution is concentrated on orbits with $R_{\rm eff}\gg 1$.
\item The $r^2$ bias asymmetry \eqref{eq:rsq-asymmetry} between $\sigma_w$ and
$\sigma_z$ mismatches survives.
\end{itemize}

\subsection{When the harmonic scaling breaks: the $\sigma_z\sim h_{\rm dm}$ regime}

The scaling fails when $R_{\rm eff}$ becomes order unity. This happens precisely when
the cloud samples large-amplitude orbits whose frequency $\Omega(J)$ is much smaller
than the small-amplitude harmonic limit. Specifically, when $\sigma_z$ approaches the
disk scale-height (or the DM scale $h_{\rm dm}$), the action distribution moves to low
$\Omega(J)$, $R_{\rm eff}$ drops, and the $r^2$ amplification of the $z$-stiffness
becomes weaker. \textbf{This is the regime where you observed $\sigma_z$ mismatch
starting to bias the fit} (the $\sigma_z=250\,\mathrm{pc}$ test): the previously
overwhelmed $\sigma_z$ direction became a real lever.

\subsection{Numerical recipe to verify in the real potential}

You can verify the $R_{\rm eff}^2$ stiffness asymmetry without trusting any
approximation, using only the simulation:
\begin{enumerate}
\item Fix $h_t$ and generate the data (one cloud).
\item Pick a fine grid $\{h_k\}$ around $h_t$. For each $h_k$, back-integrate the
data (no fitting) and compute $V_z(h_k)$ and $V_w(h_k)$.
\item Fit a parabola in $h$ to each curve; read off $V_z''(h_t)$ and $V_w''(h_t)$.
\item Compute the per-$\sigma^2$ ratio
\begin{equation}
\rho \;\equiv\; \frac{V_z''(h_t)/\sigma_z^2}{|V_w''(h_t)|/\sigma_w^2}.
\end{equation}
\item Compute $R_{\rm eff}^2$ directly: from the data cloud, evaluate the energy of
each star, find its $\Omega(J;h_t)$ (a 1D quadrature in the potential), and form
$R_{\rm eff}^2 = \sigma_w^2 / (\sigma_z^2\,\langle\Omega^2\rangle)$.
\item If $\rho \approx R_{\rm eff}^2$, the harmonic scaling is confirmed in the real
potential.
\end{enumerate}

% =========================================================
\section{Summary}
% =========================================================

\begin{enumerate}
\item The MLE balance \eqref{eq:gradbal} is exact in any potential.
\item In the harmonic case, $V_z(\nu)$ and $V_w(\nu)$ are computed in closed form
\eqref{eq:zw-back-decomp}. First derivatives \eqref{eq:firstderiv} have equal
magnitude and opposite sign — the truth-sigmas balance them. Second derivatives
satisfy the stiffness asymmetry \eqref{eq:Rsquared}: $V_z$ is $R^2$-times stiffer
per unit prior weight than $V_w$.
\item Mismatch by factor $r$ produces a numerator (force) that is symmetric between
$\sigma_z$ and $\sigma_w$ ($\propto r^2-1$), but a denominator (curvature) amplified
by $r^2$ when $\sigma_z$ is mismatched (stiff direction) and barely changed when
$\sigma_w$ is mismatched (soft direction). Hence the $r^2$ bias asymmetry
\eqref{eq:rsq-asymmetry}.
\item In the real sech$^2$+DM potential, the same structure follows from the
variational equation \eqref{eq:vareq}, with $\nu\to\Omega(J)$ averaged over the
action distribution. The harmonic $R$ becomes an orbit-averaged $R_{\rm eff}$, but
the $r^2$ bias asymmetry survives whenever the cloud sits on orbits with
$R_{\rm eff}\gg 1$.
\item The asymmetry breaks down when $\sigma_z$ becomes large compared to the disk
scale-height (the cloud reaches low-$\Omega$ orbits) — exactly the regime where
your $\sigma_z=250\,\mathrm{pc}$ test showed bias re-emerging.
\end{enumerate}

\end{document}
